\documentclass[11pt]{article}
\usepackage[margin=1in]{geometry}
\usepackage{graphicx}
\usepackage{newtxtext}
\usepackage{hyperref}
\hypersetup{colorlinks=true, urlcolor=blue}
\setlength{\parindent}{0pt}
\setlength{\parskip}{0.65\baselineskip}
\raggedright
\pagestyle{empty}

\begin{document}

\noindent
\begin{minipage}[c]{0.60\textwidth}
\includegraphics[width=\linewidth]{hkust_econ_header.png}
\end{minipage}%
\hfill
\begin{minipage}[c]{0.36\textwidth}
\raggedleft
\setlength{\parskip}{0pt}
{\bfseries Ke Cheng Hsu}\\[0.15em]
Ph.D.\ Candidate\\
LSK 5067, HKUST\\
\href{mailto:khsu@connect.ust.hk}{khsu@connect.ust.hk}\\
\href{https://cliohsu1997.github.io/}{cliohsu1997.github.io}\\[0.9em]
Oct 16, 2026
\end{minipage}

\vspace{1.1em}
{\setlength{\parskip}{0.15em}
School of Economics\\
Zhejiang University\\
No.\ 38 Zheda Road, Hangzhou, Zhejiang, China
}

To whom it may concern:

I am writing to apply for the Assistant Professor position in the School of Economics at Zhejiang University. I work in industrial organization and digital economics, with a primary interest in e-commerce return policy. I expect to complete my Ph.D.\ in Economics at the Hong Kong University of Science and Technology in 2027.

In my job market paper, ``Designing Return Policy in E-Commerce: A Trade-off between Purchase and Return Efficiency,'' I study the effect of return policies by asking whether they encourage efficient returns and purchases. A stricter return policy curbs inefficient ex-post returns, but the resulting lower return rates make higher product prices more attractive for sellers, leading to possible exclusion of ex-ante efficient purchases. Using transaction data from a mid-sized women's clothing retailer on Taobao from 2017 to 2023, I estimate a structural model of purchase, insurance, and return under monopoly pricing across six apparel categories. A lenient return policy performs better when margins and salvage value are high, by expanding valuable purchases, whereas a strict return policy performs better when margins are thinner, by curbing excessive returns. A menu that pairs a discounted price with a restocking fee balances purchase and return efficiency and outperforms both for most categories.

I am currently working on three more projects: ``Too Much Unsold Stock?: Environmental Consequences of Price Commitment'' (with Kohei Kawaguchi and Kei Kawai), ``Product Design with Heterogeneous Search Costs'' (with Pak Hung AU), and ``Pass-Through and Incidence with Product Returns'' (with Ce Bian and Liang Zhao). My future research agenda is to study how return policy should be designed when search, product substitution, and the two-sided network effect are taken into account.

I have experience teaching undergraduate macroeconomics, the economics of uncertainty and information, and Ph.D.\ microeconomic theory, and I am eager to teach microeconomic theory and industrial organization from the undergraduate to the Ph.D.\ level, and information economics through the master's level.

I will attend the China International Conference in Economics on December 12--13, 2026, and the Econometric Society Australasia Meeting on November 24--27, 2026. I will be available for interviews. I can be reached at khsu@connect.ust.hk. I look forward to your response and appreciate your consideration.

Sincerely,\\[1.4em]
Ke Cheng Hsu

\end{document}
